\section{Paměťová hierarchie}\label{sec:hw-pametova-hierarchie}

Žádný běžný druh paměti současně nenabízí nejvyšší rychlost, obrovskou kapacitu, nízkou cenu a~trvalé uložení dat. Počítač proto kombinuje několik paměťových úrovní. Blízko procesoru se nacházejí malé a~rychlé registry a~cache, pod nimi větší hlavní paměť a~ještě níže dlouhodobá úložiště. Pro zálohování a~archivaci lze přidat síťová úložiště, magnetické pásky nebo jiné systémy s~velkou kapacitou, u~nichž nevadí delší doba přístupu.

\begin{figure}[ht]
    \centering
    \begin{tikzpicture}[
    x=1cm,
    y=1cm,
    >=Stealth,
    every node/.style={font=\small},
    level/.style={
        draw,
        thick,
        rounded corners,
        align=center,
        minimum height=0.72cm
    }
]
    \node[font=\small\bfseries] at (0,3.25) {nejblíže procesoru};

    \node[level, fill=red!15, minimum width=3.0cm] (registers) at (0,2.45)
        {\textbf{registry}};
    \node[level, fill=orange!20, minimum width=4.4cm] (cache) at (0,1.45)
        {\textbf{cache}};
    \node[level, fill=yellow!24, minimum width=6.0cm] (ram) at (0,0.45)
        {\textbf{RAM}};
    \node[level, fill=green!17, minimum width=7.7cm] (storage) at (0,-0.55)
        {\textbf{SSD / HDD}};
    \node[level, fill=blue!12, minimum width=9.4cm] (archive) at (0,-1.55)
        {\textbf{archiv / záloha / síťové úložiště}};

    \draw[->, thick] (registers.south) -- (cache.north);
    \draw[->, thick] (cache.south) -- (ram.north);
    \draw[->, thick] (ram.south) -- (storage.north);
    \draw[->, thick] (storage.south) -- (archive.north);

    \draw[->, very thick] (-6.1,-2.0) -- (-6.1,2.9);
    \node[align=center, anchor=east] at (-6.35,0.45)
        {rychlost\\ cena za bit};

    \draw[->, very thick] (6.1,2.9) -- (6.1,-2.0);
    \node[align=center, anchor=west] at (6.35,0.45)
        {kapacita\\ doba uchování};
\end{tikzpicture}

    \caption{Zjednodušená paměťová hierarchie počítače.}
    \label{fig:hw-pametova-hierarchie}
\end{figure}

Směrem vzhůru v~hierarchii na obrázku~\ref{fig:hw-pametova-hierarchie} roste rychlost a~obvykle také cena za uložený bit, zatímco kapacita klesá. Směrem dolů získáváme větší a~levnější úložný prostor, ale na data čekáme déle. Jednotlivé úrovně proto plní rozdílné role:
\begin{itemize}
    \item \textbf{registry} uchovávají hodnoty, s~nimiž právě pracují instrukce;
    \item \textbf{cache} uchovává kopie nedávno nebo blízce používaných částí paměti;
    \item \textbf{RAM} obsahuje běžící programy a~jejich pracovní data;
    \item \textbf{SSD a~HDD} dlouhodobě uchovávají programy a~soubory;
    \item \textbf{archivní a~záložní systémy} upřednostňují kapacitu, cenu a~bezpečné uchování před okamžitým přístupem.
\end{itemize}

\begin{table}[ht]
    \centering
    \begin{tabular}{|l|c|c|c|}
        \hline
        \textbf{Úroveň} & \textbf{Relativní rychlost} & \textbf{Relativní kapacita} & \textbf{Volatilní} \\ \hline
        Registry & nejvyšší & nejmenší & ano \\ \hline
        Cache & velmi vysoká & malá & ano \\ \hline
        RAM & střední & větší & ano \\ \hline
        SSD/HDD & nižší & velká & ne \\ \hline
        Archiv/záloha & nejnižší & velmi velká & ne \\ \hline
    \end{tabular}
    \caption{Kvalitativní porovnání úrovní paměťové hierarchie.}
    \label{tab:hw-pametova-hierarchie}
\end{table}

Paměťová hierarchie funguje dobře tehdy, když se na rychlejších úrovních nacházejí právě ta data, která budou brzy potřeba. Část přesunů zajišťuje hardware automaticky, jako v~případě cache. Jiné přesuny řídí operační systém, například když načítá program z~SSD do RAM. Další provádí přímo uživatel, když přesune starší data do archivu nebo vytvoří zálohu.

Z~pohledu programu může být velká část těchto přesunů skrytá. Program pracuje s~adresami v~paměti a~procesor se pokouší obsloužit požadavky co nejrychleji. Pokud jsou data v~registru nebo cache, pokračuje výpočet rychle. Při výpadku cache je potřeba čekat na RAM; pokud se potřebná část programu nebo soubor teprve načítá z~úložiště, čekání je ještě delší. Rozdíl mezi jednotlivými úrovněmi vysvětluje, proč má způsob uspořádání dat významný vliv na skutečný výkon.

\importantbox{Paměťová hierarchie vytváří dojem velké a~rychlé paměti tím, že kombinuje malou rychlou paměť s~většími pomalejšími úrovněmi. Tento dojem je nejpřesvědčivější tehdy, když program vykazuje dobrou časovou a~prostorovou lokalitu.}